\documentclass[11pt,a4paper]{article}
\usepackage[utf8]{inputenc}
\usepackage[T1]{fontenc}
\usepackage[spanish,es-minimal,es-noshorthands]{babel}
\usepackage{lmodern}
\usepackage[sfdefault,lining]{FiraSans}
\usepackage[a4paper,top=2.3cm,bottom=2.2cm,left=1.8cm,right=1.8cm,headheight=1.3cm,headsep=0.4cm,footskip=1cm]{geometry}
\usepackage{amsmath,amssymb}
\usepackage{icomma}
\usepackage{graphicx}
\usepackage[table]{xcolor}
\usepackage{tikz}
\usepackage{fontawesome5}
\usepackage{ccicons}
\usepackage{fancyhdr}
\usepackage{lastpage}
\usepackage{tasks}
\usepackage{needspace}
\usepackage[most]{tcolorbox}
\usepackage{array}
\usepackage[hidelinks]{hyperref}
\AtBeginDocument{\hypersetup{pdftitle={Ficha de ejercicios de la unidad 5: Ecuaciones. 3º ESO},pdfsubject={Matemáticas 3º ESO},pdfkeywords={ecuaciones, identidades, ecuaciones de primer grado, denominadores, ecuaciones de segundo grado, discriminante, ecuaciones de tercer grado, regla de Ruffini, problemas}}}

\definecolor{navy}{HTML}{1B2748}
\definecolor{accent}{HTML}{E8603A}
\definecolor{teal}{HTML}{1F9E90}
\definecolor{soft}{HTML}{EEF1F8}
\definecolor{muted}{HTML}{5D6685}

\setlength{\parindent}{0pt}
\setlength{\parskip}{3pt}
\settasks{label=\alph*), label-format=\bfseries\color{navy}, label-width=1.3em, item-indent=1.8em, column-sep=0.6em, after-item-skip=5pt, before-skip=3pt, after-skip=4pt}

% ---------- capçalera i peu ----------
\newcommand{\llicencia}{\href{https://creativecommons.org/licenses/by-nc-sa/4.0/deed.es}{\ccbyncsa}\ \ Licencia Creative Commons Atribución-NoComercial-CompartirIgual 4.0 Internacional (CC BY-NC-SA 4.0)}
\pagestyle{fancy}
\fancyhf{}
\fancyhead[L]{\raisebox{-0.35\height}{\includegraphics[height=0.95cm]{logo-gva.png}}}
\fancyhead[C]{\small\color{muted}Ficha de ejercicios · Unidad 5: Ecuaciones · 3\textsuperscript{o} ESO}
\fancyhead[R]{\raisebox{-0.35\height}{\includegraphics[height=0.95cm]{logo-vermellar.png}}}
\fancyfoot[L]{\scriptsize\color{muted}\llicencia}
\fancyfoot[R]{\small\color{muted}\thepage}
\renewcommand{\headrulewidth}{0.4pt}
\renewcommand{\headrule}{\color{muted!40}\hrule height \headrulewidth}
\fancypagestyle{primera}{\fancyhf{}\renewcommand{\headrulewidth}{0pt}
  \fancyfoot[L]{\scriptsize\color{muted}\llicencia}
  \fancyfoot[R]{\small\color{muted}\thepage}}

% ---------- nivells i exercicis ----------
\newcommand{\estrelles}[1]{{\color{accent}\ifcase#1\or\faStar\faStar[regular]\faStar[regular]\or\faStar\faStar\faStar[regular]\or\faStar\faStar\faStar\fi}}
\newcounter{exer}
\newcommand{\ex}[1]{\par\needspace{4\baselineskip}\medskip\stepcounter{exer}\noindent{\color{navy}\bfseries\theexer.}\ \ {\small\estrelles{#1}}\ \ }

\newcommand{\seccio}[2]{\par\needspace{10\baselineskip}\bigskip
  \begin{tcolorbox}[colback=navy,colframe=navy,coltext=white,arc=2mm,boxrule=0pt,left=3mm,right=3mm,top=1.2mm,bottom=1.2mm]
  {\large\bfseries #1\quad #2}
  \end{tcolorbox}\vspace{-2pt}}

\newcolumntype{C}[1]{>{\centering\arraybackslash}p{#1}}

% quadre de recordatori (només a la fitxa d'exercicis)
\newcommand{\recorda}[1]{\par\smallskip\begin{tcolorbox}[colback=accent!7,colframe=accent!55,arc=2mm,boxrule=0.5pt,left=2.5mm,right=2.5mm,top=1mm,bottom=1mm]\small{\color{accent}\faLightbulb[regular]}\ \textbf{Recuerda:}\ #1\end{tcolorbox}}
\newcommand{\res}[1]{#1}
\newcommand{\rest}[1]{#1}

\def\fmneed{14}
% =====================================================================
% Capçalera FentMates de les fitxes: un color per a cada curs.
% Els colors són els de la pàgina inicial (LEVELS a assets/js/materials.js);
% el logo ve de _eines/logo/pdf/ (python3 _eines/logo/genera.py).
%
% Ús, al final del preàmbul de la unitat:
%   % =====================================================================
% Capçalera FentMates de les fitxes: un color per a cada curs.
% Els colors són els de la pàgina inicial (LEVELS a assets/js/materials.js);
% el logo ve de _eines/logo/pdf/ (python3 _eines/logo/genera.py).
%
% Ús, al final del preàmbul de la unitat:
%   % =====================================================================
% Capçalera FentMates de les fitxes: un color per a cada curs.
% Els colors són els de la pàgina inicial (LEVELS a assets/js/materials.js);
% el logo ve de _eines/logo/pdf/ (python3 _eines/logo/genera.py).
%
% Ús, al final del preàmbul de la unitat:
%   \input{../comu/fentmates.tex}
%   \fmcurs{1eso}
% i, al començament del document:
%   \capcalera{Matemàtiques}{1r ESO}{Fitxa d'exercicis · Unitat 2}{Divisibilitat}
% Opcional abans de \input: \def\fmneed{14} (espai mínim abans de cada secció).
% =====================================================================
\makeatletter
\def\fm@logodir{../../logo/pdf}
\providecommand{\fmneed}{10}
\def\fm@cursnom{}

% clau del curs, color i color del text damunt del color
\newcommand{\fm@curs}[3]{\@namedef{fm@color@#1}{#2}\@namedef{fm@text@#1}{#3}}
\fm@curs{1eso}{E6B84A}{navy}
\fm@curs{2eso}{3B82F6}{white}
\fm@curs{3eso}{FF7A50}{white}
\fm@curs{4eso-a}{A78BFA}{white}
\fm@curs{4eso-b}{6D28D9}{white}
\fm@curs{1batx-mi}{9B1C1C}{white}
\fm@curs{1batx-ccss}{EF6B6B}{white}
\fm@curs{2batx-mii}{1E6B47}{white}
\fm@curs{2batx-ccss}{2BB3A3}{white}

\newcommand{\fmcurs}[1]{%
  \@ifundefined{fm@color@#1}{\PackageError{fentmates}{Curs desconegut: #1}{}}{}%
  \edef\fm@tmp{\noexpand\definecolor{curs}{HTML}{\@nameuse{fm@color@#1}}%
               \noexpand\colorlet{curstext}{\@nameuse{fm@text@#1}}}\fm@tmp
  \colorlet{cursfosc}{curs!72!black}%
  \def\fm@logo{\fm@logodir/vertical-#1.pdf}%
  \def\fm@marca{\fm@logodir/marca-#1.pdf}}

% ---------- capçalera de la primera pàgina ----------
% #1 matèria (petit), #2 curs, #3 tipus de document i unitat, #4 títol
\newcommand{\capcalera}[4]{%
  \gdef\fm@cursnom{#2}%
  \par\noindent
  \begin{tcolorbox}[colback=curs,colframe=curs,boxrule=0pt,arc=2.2mm,width=3.5cm,height=2.15cm,
      valign=center,left=3mm,right=2mm,top=1mm,bottom=1mm,nobeforeafter,box align=center]
    \color{curstext}{\footnotesize #1}\par\vspace{1pt}{\fontsize{19}{21}\selectfont\bfseries #2\par}%
  \end{tcolorbox}\hspace{5mm}%
  \begin{minipage}[c]{\dimexpr\linewidth-3.5cm-5mm-2.9cm\relax}
    {\color{cursfosc}\small\bfseries\MakeUppercase{#3}}\par\vspace{2pt}
    {\color{navy}\fontsize{20}{23}\selectfont\bfseries #4\par}
  \end{minipage}\hfill
  \begin{minipage}[c]{2.6cm}\raggedleft\includegraphics[width=2.6cm]{\fm@logo}\end{minipage}\par
  \vspace{2pt}}

% ---------- barres de secció del color del curs ----------
\renewcommand{\seccio}[2]{\par\needspace{\fmneed\baselineskip}\bigskip
  \begin{tcolorbox}[colback=curs,colframe=curs,coltext=curstext,arc=2mm,boxrule=0pt,left=3mm,right=3mm,top=1.2mm,bottom=1.2mm]
  {\large\bfseries #1\quad #2}
  \end{tcolorbox}\vspace{-2pt}}

% ---------- estrelles de dificultat del color del curs (a Batxillerat no se'n posen) ----------
\@ifundefined{estrelles}{}{%
  \renewcommand{\estrelles}[1]{{\color{curs}\ifcase#1\or\faStar\faStar[regular]\faStar[regular]\or\faStar\faStar\faStar[regular]\or\faStar\faStar\faStar\fi}}}

% ---------- capçalera i peu de les pàgines ----------
% primera pàgina: logos institucionals al peu
\fancypagestyle{primera}{\fancyhf{}\renewcommand{\headrulewidth}{0pt}
  \fancyfoot[L]{\raisebox{-0.32\height}{\includegraphics[height=0.78cm]{logo-gva.png}}\hspace{5mm}%
    \raisebox{-0.32\height}{\includegraphics[height=0.78cm]{logo-vermellar.png}}}
  \fancyfoot[R]{\scriptsize\color{muted}\href{https://creativecommons.org/licenses/by-nc-sa/4.0/deed.ca}{\ccbyncsa}\ \ CC BY-NC-SA 4.0 · fentmates.org\qquad\small\thepage}}
% la resta: el símbol a l'esquerra i el curs a la dreta, amb el títol al centre com abans
\fancyhead[L]{\raisebox{-0.3\height}{\includegraphics[height=0.62cm]{\fm@marca}}}
\fancyhead[R]{\small\bfseries\color{cursfosc}\fm@cursnom}
\renewcommand{\headrule}{\color{curs!45}\hrule height \headrulewidth}
\makeatother

%   \fmcurs{1eso}
% i, al començament del document:
%   \capcalera{Matemàtiques}{1r ESO}{Fitxa d'exercicis · Unitat 2}{Divisibilitat}
% Opcional abans de \input: \def\fmneed{14} (espai mínim abans de cada secció).
% =====================================================================
\makeatletter
\def\fm@logodir{../../logo/pdf}
\providecommand{\fmneed}{10}
\def\fm@cursnom{}

% clau del curs, color i color del text damunt del color
\newcommand{\fm@curs}[3]{\@namedef{fm@color@#1}{#2}\@namedef{fm@text@#1}{#3}}
\fm@curs{1eso}{E6B84A}{navy}
\fm@curs{2eso}{3B82F6}{white}
\fm@curs{3eso}{FF7A50}{white}
\fm@curs{4eso-a}{A78BFA}{white}
\fm@curs{4eso-b}{6D28D9}{white}
\fm@curs{1batx-mi}{9B1C1C}{white}
\fm@curs{1batx-ccss}{EF6B6B}{white}
\fm@curs{2batx-mii}{1E6B47}{white}
\fm@curs{2batx-ccss}{2BB3A3}{white}

\newcommand{\fmcurs}[1]{%
  \@ifundefined{fm@color@#1}{\PackageError{fentmates}{Curs desconegut: #1}{}}{}%
  \edef\fm@tmp{\noexpand\definecolor{curs}{HTML}{\@nameuse{fm@color@#1}}%
               \noexpand\colorlet{curstext}{\@nameuse{fm@text@#1}}}\fm@tmp
  \colorlet{cursfosc}{curs!72!black}%
  \def\fm@logo{\fm@logodir/vertical-#1.pdf}%
  \def\fm@marca{\fm@logodir/marca-#1.pdf}}

% ---------- capçalera de la primera pàgina ----------
% #1 matèria (petit), #2 curs, #3 tipus de document i unitat, #4 títol
\newcommand{\capcalera}[4]{%
  \gdef\fm@cursnom{#2}%
  \par\noindent
  \begin{tcolorbox}[colback=curs,colframe=curs,boxrule=0pt,arc=2.2mm,width=3.5cm,height=2.15cm,
      valign=center,left=3mm,right=2mm,top=1mm,bottom=1mm,nobeforeafter,box align=center]
    \color{curstext}{\footnotesize #1}\par\vspace{1pt}{\fontsize{19}{21}\selectfont\bfseries #2\par}%
  \end{tcolorbox}\hspace{5mm}%
  \begin{minipage}[c]{\dimexpr\linewidth-3.5cm-5mm-2.9cm\relax}
    {\color{cursfosc}\small\bfseries\MakeUppercase{#3}}\par\vspace{2pt}
    {\color{navy}\fontsize{20}{23}\selectfont\bfseries #4\par}
  \end{minipage}\hfill
  \begin{minipage}[c]{2.6cm}\raggedleft\includegraphics[width=2.6cm]{\fm@logo}\end{minipage}\par
  \vspace{2pt}}

% ---------- barres de secció del color del curs ----------
\renewcommand{\seccio}[2]{\par\needspace{\fmneed\baselineskip}\bigskip
  \begin{tcolorbox}[colback=curs,colframe=curs,coltext=curstext,arc=2mm,boxrule=0pt,left=3mm,right=3mm,top=1.2mm,bottom=1.2mm]
  {\large\bfseries #1\quad #2}
  \end{tcolorbox}\vspace{-2pt}}

% ---------- estrelles de dificultat del color del curs (a Batxillerat no se'n posen) ----------
\@ifundefined{estrelles}{}{%
  \renewcommand{\estrelles}[1]{{\color{curs}\ifcase#1\or\faStar\faStar[regular]\faStar[regular]\or\faStar\faStar\faStar[regular]\or\faStar\faStar\faStar\fi}}}

% ---------- capçalera i peu de les pàgines ----------
% primera pàgina: logos institucionals al peu
\fancypagestyle{primera}{\fancyhf{}\renewcommand{\headrulewidth}{0pt}
  \fancyfoot[L]{\raisebox{-0.32\height}{\includegraphics[height=0.78cm]{logo-gva.png}}\hspace{5mm}%
    \raisebox{-0.32\height}{\includegraphics[height=0.78cm]{logo-vermellar.png}}}
  \fancyfoot[R]{\scriptsize\color{muted}\href{https://creativecommons.org/licenses/by-nc-sa/4.0/deed.ca}{\ccbyncsa}\ \ CC BY-NC-SA 4.0 · fentmates.org\qquad\small\thepage}}
% la resta: el símbol a l'esquerra i el curs a la dreta, amb el títol al centre com abans
\fancyhead[L]{\raisebox{-0.3\height}{\includegraphics[height=0.62cm]{\fm@marca}}}
\fancyhead[R]{\small\bfseries\color{cursfosc}\fm@cursnom}
\renewcommand{\headrule}{\color{curs!45}\hrule height \headrulewidth}
\makeatother

%   \fmcurs{1eso}
% i, al començament del document:
%   \capcalera{Matemàtiques}{1r ESO}{Fitxa d'exercicis · Unitat 2}{Divisibilitat}
% Opcional abans de \input: \def\fmneed{14} (espai mínim abans de cada secció).
% =====================================================================
\makeatletter
\def\fm@logodir{../../logo/pdf}
\providecommand{\fmneed}{10}
\def\fm@cursnom{}

% clau del curs, color i color del text damunt del color
\newcommand{\fm@curs}[3]{\@namedef{fm@color@#1}{#2}\@namedef{fm@text@#1}{#3}}
\fm@curs{1eso}{E6B84A}{navy}
\fm@curs{2eso}{3B82F6}{white}
\fm@curs{3eso}{FF7A50}{white}
\fm@curs{4eso-a}{A78BFA}{white}
\fm@curs{4eso-b}{6D28D9}{white}
\fm@curs{1batx-mi}{9B1C1C}{white}
\fm@curs{1batx-ccss}{EF6B6B}{white}
\fm@curs{2batx-mii}{1E6B47}{white}
\fm@curs{2batx-ccss}{2BB3A3}{white}

\newcommand{\fmcurs}[1]{%
  \@ifundefined{fm@color@#1}{\PackageError{fentmates}{Curs desconegut: #1}{}}{}%
  \edef\fm@tmp{\noexpand\definecolor{curs}{HTML}{\@nameuse{fm@color@#1}}%
               \noexpand\colorlet{curstext}{\@nameuse{fm@text@#1}}}\fm@tmp
  \colorlet{cursfosc}{curs!72!black}%
  \def\fm@logo{\fm@logodir/vertical-#1.pdf}%
  \def\fm@marca{\fm@logodir/marca-#1.pdf}}

% ---------- capçalera de la primera pàgina ----------
% #1 matèria (petit), #2 curs, #3 tipus de document i unitat, #4 títol
\newcommand{\capcalera}[4]{%
  \gdef\fm@cursnom{#2}%
  \par\noindent
  \begin{tcolorbox}[colback=curs,colframe=curs,boxrule=0pt,arc=2.2mm,width=3.5cm,height=2.15cm,
      valign=center,left=3mm,right=2mm,top=1mm,bottom=1mm,nobeforeafter,box align=center]
    \color{curstext}{\footnotesize #1}\par\vspace{1pt}{\fontsize{19}{21}\selectfont\bfseries #2\par}%
  \end{tcolorbox}\hspace{5mm}%
  \begin{minipage}[c]{\dimexpr\linewidth-3.5cm-5mm-2.9cm\relax}
    {\color{cursfosc}\small\bfseries\MakeUppercase{#3}}\par\vspace{2pt}
    {\color{navy}\fontsize{20}{23}\selectfont\bfseries #4\par}
  \end{minipage}\hfill
  \begin{minipage}[c]{2.6cm}\raggedleft\includegraphics[width=2.6cm]{\fm@logo}\end{minipage}\par
  \vspace{2pt}}

% ---------- barres de secció del color del curs ----------
\renewcommand{\seccio}[2]{\par\needspace{\fmneed\baselineskip}\bigskip
  \begin{tcolorbox}[colback=curs,colframe=curs,coltext=curstext,arc=2mm,boxrule=0pt,left=3mm,right=3mm,top=1.2mm,bottom=1.2mm]
  {\large\bfseries #1\quad #2}
  \end{tcolorbox}\vspace{-2pt}}

% ---------- estrelles de dificultat del color del curs (a Batxillerat no se'n posen) ----------
\@ifundefined{estrelles}{}{%
  \renewcommand{\estrelles}[1]{{\color{curs}\ifcase#1\or\faStar\faStar[regular]\faStar[regular]\or\faStar\faStar\faStar[regular]\or\faStar\faStar\faStar\fi}}}

% ---------- capçalera i peu de les pàgines ----------
% primera pàgina: logos institucionals al peu
\fancypagestyle{primera}{\fancyhf{}\renewcommand{\headrulewidth}{0pt}
  \fancyfoot[L]{\raisebox{-0.32\height}{\includegraphics[height=0.78cm]{logo-gva.png}}\hspace{5mm}%
    \raisebox{-0.32\height}{\includegraphics[height=0.78cm]{logo-vermellar.png}}}
  \fancyfoot[R]{\scriptsize\color{muted}\href{https://creativecommons.org/licenses/by-nc-sa/4.0/deed.ca}{\ccbyncsa}\ \ CC BY-NC-SA 4.0 · fentmates.org\qquad\small\thepage}}
% la resta: el símbol a l'esquerra i el curs a la dreta, amb el títol al centre com abans
\fancyhead[L]{\raisebox{-0.3\height}{\includegraphics[height=0.62cm]{\fm@marca}}}
\fancyhead[R]{\small\bfseries\color{cursfosc}\fm@cursnom}
\renewcommand{\headrule}{\color{curs!45}\hrule height \headrulewidth}
\makeatother

\fmcurs{3eso}

\begin{document}
\thispagestyle{primera}

\capcalera{Matemáticas}{3\textsuperscript{o} ESO}{Ficha de ejercicios · Unidad 5}{Ecuaciones}
\vspace{10pt}
Nombre y apellidos: \rule{6.6cm}{0.4pt}\hfill Grupo: \rule{1.6cm}{0.4pt}\hfill Fecha: \rule{2.2cm}{0.4pt}
\vspace{8pt}

\begin{tcolorbox}[colback=soft,colframe=soft,arc=2mm,left=3mm,right=3mm,top=1.5mm,bottom=1.5mm]
\small Ejercicios graduados por niveles de dificultad:\quad \estrelles{1}\ básico \qquad \estrelles{2}\ medio \qquad \estrelles{3}\ avanzado.\\[2pt] Resuélvelos en el cuaderno mostrando todos los pasos y comprueba las soluciones siempre que puedas. Haz los cálculos sin calculadora.
\end{tcolorbox}

% =====================================================
\seccio{1}{Ecuaciones e identidades}

\ex{1} Indica qué igualdades son identidades y cuáles son ecuaciones:\begin{tasks}(2)
  \task $3(x+2)=3x+6$
  \task $3x+2=8$
  \task $(x+1)^2=x^2+2x+1$
  \task $x^2-1=0$
  \task $5x-2x=3x$
  \task $5x-2=3x$
  \task $(x-3)(x+3)=x^2-9$
  \task $2(x-1)=2x-1$
\end{tasks}\recorda{Una \textbf{identidad} se cumple para cualquier valor de $x$; una \textbf{ecuación}, solo para algunos valores (o para ninguno).}

\ex{1} Comprueba si los valores indicados son solución de cada ecuación:\begin{tasks}(2)
  \task $3x-5=x+1$, \ $x=3$
  \task $x^2-4x=5$, \ $x=-1$ y $x=5$
  \task $2(x-1)=x+4$, \ $x=5$
  \task $x^3-2x=4$, \ $x=2$
\end{tasks}

\ex{1} En la ecuación $5x^2-3x+1=2x-7$, indica el primer y el segundo miembro, los términos, la incógnita y el grado.

\ex{2} Indica el grado de cada ecuación una vez reducida:\begin{tasks}(2)
  \task $x(x+4)=x^2-8$
  \task $(x-2)^2=x^2+4$
  \task $x^3-x(x^2+1)=2$
  \task $(x+1)(x-1)=3x$
\end{tasks}\recorda{Antes de decir el grado, opera y reduce: los términos de grado alto se pueden anular.}

\ex{2} ¿Son equivalentes estas parejas de ecuaciones? Justifícalo.\begin{tasks}(1)
  \task $2x-3=7$ \ y \ $3x+1=16$
  \task $x^2=25$ \ y \ $x=5$
\end{tasks}

\ex{2} Escribe una ecuación de primer grado que tenga la solución $x=-2$ y una ecuación de segundo grado que tenga las soluciones $x=3$ y $x=-3$.

\ex{3} Halla el valor de $k$ en cada caso:\begin{tasks}(1)
  \task $x=3$ es solución de $2x^2-kx+3=0$.
  \task $x=-1$ es solución de $kx^3+2x=k-6$.
\end{tasks}


% =====================================================
\seccio{2}{Ecuaciones de primer grado}

\ex{1} Resuelve:\begin{tasks}(4)
  \task $x+7=12$
  \task $3x=-21$
  \task $5x-4=11$
  \task $-2x+9=1$
  \task $4x+3=2x+11$
  \task $7-3x=x-13$
  \task $6x-5=9x+7$
  \task $2x-1=5x-1$
\end{tasks}\recorda{Lo que suma pasa restando y lo que resta pasa sumando; lo que multiplica a todo un miembro pasa dividiendo (con su signo).}

\ex{1} Resuelve quitando primero los paréntesis:\begin{tasks}(2)
  \task $3(x-2)=12$
  \task $5(x+1)=2x-7$
  \task $4-(x+3)=2x-5$
  \task $2(3x-1)=4(x+2)$
\end{tasks}

\ex{2} Resuelve:\begin{tasks}(2)
  \task $3(2x-1)-2(x+4)=5x-13$
  \task $5-2(3-x)=4(x+1)-11$
  \task $x-[2x-(3-x)]=7$
  \task $4(x-3)-(x-5)=2(x+1)$
\end{tasks}\recorda{Un signo $-$ delante de un paréntesis cambia el signo de \textbf{todos} los términos de dentro.}

\ex{2} Resuelve y di qué ecuaciones no tienen solución y cuáles tienen infinitas:\begin{tasks}(2)
  \task $3(x+2)-x=2(x+3)$
  \task $5x-2(x-1)=3x+4$
  \task $6x-2(3x+1)=-2$
  \task $7x-3(2x-1)=2x+1$
\end{tasks}

\ex{2} Resuelve y comprueba la solución:\begin{tasks}(2)
  \task $(x+1)^{2}-(x-1)^{2}=12$
  \task $x(x+3)-x^{2}=2x+5$
\end{tasks}

\ex{3} Resuelve:\begin{tasks}(1)
  \task $2[3-2(x+1)]=5-3(2x-1)-x$
  \task $(x-3)(x+3)-(x-1)^{2}=2(x+1)$
\end{tasks}

\ex{3} Halla el valor de $a$ para que la ecuación $3(x-a)=2x+1$ tenga la solución $x=5$.


% =====================================================
\seccio{3}{Ecuaciones de primer grado con denominadores}

\ex{1} Resuelve:\begin{tasks}(4)
  \task $\dfrac{x}{3}=4$
  \task $\dfrac{x}{2}+\dfrac{x}{4}=6$
  \task $x-\dfrac{x}{3}=10$
  \task $\dfrac{2x}{5}-1=3$
\end{tasks}

\ex{1} Resuelve:\begin{tasks}(4)
  \task $\dfrac{x+2}{3}=4$
  \task $\dfrac{x-1}{2}=\dfrac{x+3}{4}$
  \task $\dfrac{3x-1}{4}=2$
  \task $\dfrac{x}{6}+1=\dfrac{x}{4}$
\end{tasks}\recorda{Multiplica \textbf{todos} los términos por el m.c.m. de los denominadores (también los que no tienen denominador) y pon los numeradores entre paréntesis.}

\ex{2} Resuelve:\begin{tasks}(2)
  \task $\dfrac{x+1}{2}-\dfrac{x-1}{3}=2$
  \task $\dfrac{x-2}{4}-\dfrac{x-3}{6}=\dfrac{1}{2}$
  \task $\dfrac{2x+1}{3}-\dfrac{x-2}{2}=3$
  \task $x-\dfrac{x+2}{5}=\dfrac{x+6}{2}$
\end{tasks}

\ex{2} Resuelve:\begin{tasks}(2)
  \task $\dfrac{x-1}{3}+\dfrac{x+2}{6}=x-2$
  \task $3-\dfrac{x-1}{4}=\dfrac{x+5}{2}$
\end{tasks}

\ex{3} Resuelve e interpreta el resultado:\begin{tasks}(1)
  \task $\dfrac{3(x-2)}{4}-\dfrac{2(x+1)}{3}=\dfrac{x-10}{6}$
  \task $\dfrac{x}{2}-\dfrac{x-1}{3}-\dfrac{x+1}{4}=\dfrac{1-x}{12}$
  \task $\dfrac{x+1}{2}-\dfrac{x+3}{4}=\dfrac{x}{4}-1$
\end{tasks}

\ex{3} Un alumno ha resuelto la ecuación $\dfrac{x-2}{3}-\dfrac{x+1}{2}=1$ así:\quad $2x-4-3x+3=6\ \Rightarrow\ -x=7\ \Rightarrow\ x=-7$. \ ¿Dónde se ha equivocado? Resuélvela correctamente.


% =====================================================
\seccio{4}{Ecuaciones de segundo grado incompletas}

\ex{1} Escribe cada ecuación en la forma general $ax^2+bx+c=0$, indica los coeficientes $a$, $b$ y $c$ y di si es completa o incompleta:\begin{tasks}(2)
  \task $5x-x^2=6$
  \task $3x^2=12$
  \task $x(x-4)=0$
  \task $(x+1)^2=2x+5$
\end{tasks}

\ex{1} Resuelve:\begin{tasks}(3)
  \task $x^{2}-36=0$
  \task $2x^{2}-50=0$
  \task $x^{2}+9=0$
  \task $4x^{2}-9=0$
  \task $3x^{2}=27$
  \task $x^{2}-7=0$
\end{tasks}\recorda{En $ax^2+c=0$ despejamos $x^2$ y hacemos la raíz: hay \textbf{dos} soluciones, $x=\pm\sqrt{\ }$, o ninguna si el número es negativo.}

\ex{1} Resuelve:\begin{tasks}(4)
  \task $x^{2}-5x=0$
  \task $3x^{2}+12x=0$
  \task $2x^{2}=7x$
  \task $-x^{2}+6x=0$
\end{tasks}\recorda{En $ax^2+bx=0$ sacamos factor común $x$: una solución siempre es $x=0$. ¡No dividas entre $x$!}

\ex{2} Reduce y resuelve:\begin{tasks}(2)
  \task $(x+3)(x-3)=7$
  \task $x(x-2)=4-2x$
  \task $(2x-1)^{2}=1$
  \task $(x+2)^{2}-4=3x$
\end{tasks}

\ex{2} Resuelve:\begin{tasks}(2)
  \task $\dfrac{x^{2}}{3}=12$
  \task $\dfrac{x^{2}}{4}-\dfrac{x}{2}=0$
  \task $\dfrac{x^{2}+1}{5}=2$
  \task $\dfrac{x^{2}-2}{2}-\dfrac{x^{2}+1}{3}=1$
\end{tasks}

\ex{3} Resuelve:\begin{tasks}(1)
  \task $(2x+3)(2x-3)-(x+1)^{2}=2(x^{2}-x)-10$
  \task $\dfrac{x(x-3)}{2}-\dfrac{x(x+1)}{3}=x$
\end{tasks}

\ex{3} Marta resuelve $5x^2=3x$ dividiendo los dos miembros entre $x$ y obtiene $5x=3$, es decir, $x=\tfrac35$. ¿Qué ha hecho mal? Resuélvela correctamente.

\ex{3} Escribe una ecuación de segundo grado incompleta que tenga por soluciones:\begin{tasks}(2)
  \task $4$ y $-4$
  \task $0$ y $3$
\end{tasks}


% =====================================================
\seccio{5}{Ecuaciones de segundo grado completas}

\ex{1} Resuelve con la fórmula:\begin{tasks}(4)
  \task $x^{2}-6x+8=0$
  \task $x^{2}+2x-8=0$
  \task $x^{2}-4x-5=0$
  \task $x^{2}+7x+12=0$
\end{tasks}\recorda{$x=\dfrac{-b\pm\sqrt{b^2-4ac}}{2a}$. \ Cuidado con los signos: si $b=-6$, entonces $-b=6$ y $b^2=36$.}

\ex{1} Resuelve:\begin{tasks}(4)
  \task $2x^{2}-5x+2=0$
  \task $3x^{2}+x-2=0$
  \task $6x^{2}-x-1=0$
  \task $4x^{2}+4x-3=0$
\end{tasks}

\ex{1} Calcula el discriminante y di cuántas soluciones tiene cada ecuación, sin resolverla:\begin{tasks}(4)
  \task $x^{2}-10x+25=0$
  \task $2x^{2}+3x+4=0$
  \task $x^{2}-x-1=0$
  \task $-3x^{2}+2x+1=0$
\end{tasks}

\ex{2} Resuelve:\begin{tasks}(4)
  \task $-x^{2}+x+6=0$
  \task $9x^{2}-6x+1=0$
  \task $x^{2}+3x+5=0$
  \task $x^{2}-2x-2=0$
\end{tasks}

\ex{2} Reduce a la forma general y resuelve:\begin{tasks}(2)
  \task $(x-1)(x+4)=6$
  \task $(x+2)^{2}=3x+10$
  \task $x(2x-1)=(x+1)^{2}-4$
  \task $3x(x-2)-(x-1)(x+1)=9$
\end{tasks}

\ex{2} Resuelve:\begin{tasks}(2)
  \task $\dfrac{x^{2}}{2}-\dfrac{x}{3}=\dfrac{4}{3}$
  \task $\dfrac{x^{2}-1}{3}+\dfrac{x}{2}=2$
\end{tasks}

\ex{3} Resuelve:\begin{tasks}(1)
  \task $\dfrac{(x-1)^{2}}{4}-\dfrac{x+2}{6}=\dfrac{x-3}{3}$
  \task $\dfrac{x(x+1)}{2}-\dfrac{(x-1)(x+2)}{3}=1$
\end{tasks}

\ex{3} Responde:\begin{tasks}(1)
  \task Halla $k$ para que la ecuación $x^2-4x+k=0$ tenga una única solución. ¿Cuál es?
  \task Halla $k$ para que $x=3$ sea solución de $x^2+kx-6=0$ y calcula la otra solución.
\end{tasks}

\ex{3} Escribe una ecuación de segundo grado, en forma general, que tenga por soluciones:\begin{tasks}(2)
  \task $2$ y $-5$
  \task $\tfrac12$ y $3$
\end{tasks}


% =====================================================
\seccio{6}{Ecuaciones de tercer grado (Ruffini)}

\ex{1} Resuelve, sin desarrollar los productos:\begin{tasks}(2)
  \task $(x-2)(x+3)(x-5)=0$
  \task $x(x+1)(2x-3)=0$
  \task $(x-4)(x^2-9)=0$
  \task $(x+1)(x^2+4)=0$
\end{tasks}\recorda{Si un producto vale 0, alguno de los factores tiene que valer 0.}

\ex{1} Resuelve sacando factor común:\begin{tasks}(2)
  \task $x^{3}-16x=0$
  \task $x^{3}+3x^{2}-10x=0$
  \task $2x^{3}-8x^{2}=0$
  \task $x^{3}+x^{2}+x=0$
\end{tasks}

\ex{2} Resuelve con la regla de Ruffini:\begin{tasks}(2)
  \task $x^{3}-7x+6=0$
  \task $x^{3}-4x^{2}+x+6=0$
  \task $x^{3}+2x^{2}-5x-6=0$
  \task $x^{3}-3x^{2}-4x+12=0$
\end{tasks}\recorda{Las soluciones enteras son divisores del término independiente. Cuando encuentres una, $x=a$, divide por $x-a$ con Ruffini y resuelve la ecuación de segundo grado que queda.}

\ex{2} Resuelve. Atención a las soluciones dobles y a las ecuaciones de segundo grado sin solución:\begin{tasks}(2)
  \task $x^{3}-x^{2}-8x+12=0$
  \task $x^{3}+3x^{2}+3x+1=0$
  \task $x^{3}-2x^{2}+x-2=0$
  \task $x^{3}+x^{2}-2=0$
\end{tasks}

\ex{2} Pásalo todo al primer miembro y resuelve:\begin{tasks}(2)
  \task $x^{3}+12=4x^{2}+3x$
  \task $x(x-1)(x+1)=24$
\end{tasks}

\ex{3} Resuelve (alguna solución es una fracción):\begin{tasks}(2)
  \task $2x^{3}-5x^{2}-4x+3=0$
  \task $3x^{3}+4x^{2}-5x-2=0$
\end{tasks}

\ex{3} Aplica Ruffini dos veces para resolver estas ecuaciones de cuarto grado:\begin{tasks}(2)
  \task $x^{4}-10x^{2}+9=0$
  \task $x^{4}-x^{3}-7x^{2}+x+6=0$
\end{tasks}

\ex{3} Halla $k$ para que $x=2$ sea solución de $x^3-kx^2+2x+4=0$ y, con este valor, halla las otras soluciones.


% =====================================================
\seccio{7}{Problemas con ecuaciones}

\ex{1} La suma de dos números consecutivos es 75. ¿Cuáles son?\recorda{Elige la incógnita y escribe qué representa, plantea la ecuación, resuélvela y comprueba que la solución tiene sentido en el problema.}

\ex{1} Si a un número le sumamos su triple, obtenemos 56. ¿Cuál es el número?

\ex{1} Una persona gasta la tercera parte de su dinero en ropa y la cuarta parte en libros, y todavía le quedan 25 €. ¿Cuánto dinero tenía?

\ex{1} Ana tiene 12 años y su padre, 40. ¿Dentro de cuántos años la edad del padre será el doble que la de Ana?

\ex{2} El largo de un rectángulo es 3 m más que el doble del ancho y el perímetro mide 54 m. ¿Qué dimensiones tiene?

\ex{2} En un corral hay gallinas y conejos. En total hay 40 cabezas y 110 patas. ¿Cuántos animales hay de cada tipo?

\ex{2} En una hucha hay 30 monedas, unas de 1 € y otras de 2 €, que suman 46 €. ¿Cuántas monedas hay de cada tipo?

\ex{2} La edad de una madre es el triple de la de su hijo, y dentro de 12 años solo será el doble. ¿Cuántos años tienen ahora?

\ex{2} Halla dos números naturales consecutivos cuyo producto sea 156.

\ex{2} Si al cuadrado de un número le restamos su triple, obtenemos 40. ¿Cuál puede ser el número?

\ex{2} Si aumentamos en 3 cm el lado de un cuadrado, su área aumenta 51 cm². ¿Cuánto mide el lado?

\ex{3} La hipotenusa de un triángulo rectángulo mide 10 cm y un cateto mide 2 cm más que el otro. ¿Cuánto miden los catetos?

\ex{3} Un rectángulo tiene 34 m de perímetro y 60 m² de área. ¿Qué dimensiones tiene?

\ex{3} Queremos 20 kg de una mezcla de café a 9,50 €/kg mezclando café de 12 €/kg con café de 8 €/kg. ¿Cuántos kilos de cada tipo tenemos que poner?

\ex{3} Dos ciclistas salen a la vez de dos ciudades separadas 90 km y van el uno hacia el otro a 20 km/h y a 25 km/h. ¿Cuánto tiempo tardan en encontrarse? ¿A qué distancia de la primera ciudad?

\ex{3} Una caja de base cuadrada tiene una altura 3 cm mayor que el lado de la base, y su volumen es 112 cm³. ¿Qué dimensiones tiene?

\ex{3} El producto de tres números naturales consecutivos es 8 veces su suma. ¿Cuáles son?

\vfill
\begin{center}\color{muted}\small\faCheckCircle[regular]\ \ ¡Ánimo! Si te cuesta algún ejercicio, repasa los apuntes interactivos de la unidad.\end{center}

\end{document}
